\chapterhead{44}{ORR-5}{Risk Mitigation}{1}{2}

\lohead{44}{a}{Explain different ways firms address their operational risk
exposures.}

\orsub{Risk mitigation}

Risk mitigation involves measures to minimise operational risk to acceptable
levels, differentiating it from credit, market, and liquidity risks. There are two
main types of operational risk: \textbf{external} (out of the firm's control) and
\textbf{internal} (people, systems, processes).

\orsubb{To reduce operational risk, firms can}

\begin{center}
\begin{tikzpicture}[font=\fontsize{8.4}{10.5}\selectfont,
  bx/.style={text width=3.3cm,align=center,rounded corners=3pt,inner sep=8pt,
             text=white,font=\sffamily\bfseries\fontsize{8.6}{10.4}\selectfont}]
\node[bx,fill=navy]   at (-5.4,0) {a) Set up strong internal controls};
\node[bx,fill=teal]   at (-1.8,0) {b) Use automation};
\node[bx,fill=forest] at ( 1.8,0) {c) Buy insurance};
\node[bx,fill=olive]  at ( 5.4,0) {d) Maintain capital reserves};
\end{tikzpicture}
\end{center}

\begin{center}
\begin{tikzpicture}[font=\fontsize{8.4}{10.5}\selectfont,
  t/.style={text width=2.5cm,align=center,rounded corners=3pt,inner sep=6pt,
            text=white,font=\sffamily\bfseries\fontsize{8.6}{10.4}\selectfont}]
\node[fill=navy,text=white,rounded corners=3pt,inner sep=6pt,
      font=\sffamily\bfseries\fontsize{9}{11}\selectfont] (top) at (0,1.6)
      {Risk Responses --- the Four T's};
\node[t,fill=rust]   (a) at (-5.1,0) {Tolerate};
\node[t,fill=teal]   (b) at (-1.7,0) {Treat};
\node[t,fill=forest] (c) at ( 1.7,0) {Transfer};
\node[t,fill=olive]  (d) at ( 5.1,0) {Terminate};
\foreach \n in {a,b,c,d}
  \draw[palenavy,line width=1pt] (top.south) -- ($(\n.north)+(0,0.02)$);
\end{tikzpicture}
\end{center}
\orfigcap{The four generic responses to any identified risk.}

% ---------------------------------------------------------------------
\lohead{44}{b}{Describe and provide examples of different types of internal
controls, and explain the process of internal control design and control testing.}

\orsub{Types of internal controls}

\noindent\begin{minipage}{\textwidth}
\renewcommand{\arraystretch}{1.35}
\begin{tabularx}{\textwidth}{@{}L{0.17\textwidth}YL{0.34\textwidth}@{}}
\rowcolor{navy} \thd{Control} & \thd{What it does} & \thd{Example}\\
{\tabfont \textbf{1) Preventive}} & {\tabfont Stop problems before they happen} & {\tabfont Segregating job duties to prevent fraud}\\
\rowcolor{paleblue}
{\tabfont \textbf{2) Detective}} & {\tabfont Identify issues quickly to minimise damage} & {\tabfont Monitoring accounts for suspicious activity}\\
{\tabfont \textbf{3) Corrective}} & {\tabfont Fix problems and prevent them from recurring} & {\tabfont Data backups to recover from losses}\\
\rowcolor{paleblue}
{\tabfont \textbf{4) Directive}} & {\tabfont Guide employees on how to perform tasks correctly} & {\tabfont Training and procedures}\\
\end{tabularx}
\end{minipage}
\orfigcap{The four types of internal control.}

% ---------------------------------------------------------------------
\lohead{44}{c}{Describe control automation, internal control design, and control
testing, including risks and challenges that arise in these processes and ways to
make them more effective.}

\begin{orkeybox}
\textbf{Key controls} are especially important because they can function
independently to address risks. They can be preventive (like automation to reduce
errors) or corrective (like backups to restore lost data).\par
\vspace{5pt}
\textbf{Automated controls} are often better than manual ones due to their
reliability. However, they introduce new risks like system errors or bad data
leading to faulty controls.
\end{orkeybox}

\orsub{Control testing}

Control testing verifies if controls are designed well and used properly. Here are
some poorly designed controls to avoid in a company:

\begin{lettered}
\item \textbf{Optimistic controls:} very reliant on individuals (e.g., quick
sign-offs without proper review).
\item \textbf{Collective controls:} divide tasks but reduce accountability (e.g.,
double-checking without proper attention).
\item \textbf{More of the same approach:} adding more of a faulty control won't
solve the problem.
\end{lettered}

\orsubb{Four main categories of control testing}

\begin{lettered}
\item \textbf{Self-assessment:} suitable for low-risk situations or secondary
controls.
\item \textbf{Examination:} reviews documentation and testing results. Best for
automated controls.
\item \textbf{Observation:} witness testing of controls in action. Best for key
controls.
\item \textbf{Reperformance:} testing controls on a sample of transactions (e.g.,
introducing errors to see if they're detected).
\end{lettered}

\noindent\begin{minipage}{\textwidth}
\renewcommand{\arraystretch}{1.32}
\begin{tabularx}{\textwidth}{@{}YYY@{}}
\rowcolor{navy} \thd{Who does the testing?} & \thd{How often?} & \thd{Three lines of defense roles}\\
{\tabfont
a)~Generally, an independent person (not the control designer), to avoid
bias.\newline
b)~Self-certification is okay sometimes, but rare.} &
{\tabfont
a)~More frequent for higher risk, and vice versa.} &
{\tabfont
a)~First line can test key controls they use every day.\newline
b)~ORM (2nd line) gets results, not involved in testing itself.\newline
c)~Internal Audit (3rd line) usually tests controls.}\\
\end{tabularx}
\end{minipage}
\orfigcap{Who tests controls, how often, and where responsibility sits.}

% ---------------------------------------------------------------------
\lohead{44}{d}{Describe methods to improve the quality of an operational process
and reduce the potential for human error.}

\orsub{Process design and categories of human errors}

\begin{ordefbox}{Prevention through Design (PtD)}
Optimise processes to minimise errors through checklists, protocols, systems, and
standardisation. This tackles human error at the source, similar to the Swiss
Cheese Model: multiple safeguards, each with holes.
\end{ordefbox}

\orsubb{Categorising human errors}

\noindent\begin{minipage}{\textwidth}
\renewcommand{\arraystretch}{1.35}
\begin{tabularx}{\textwidth}{@{}L{0.24\textwidth}YL{0.32\textwidth}@{}}
\rowcolor{navy} \thd{Error type} & \thd{What it is} & \thd{Remedies}\\
{\tabfont \textbf{1) Slips}} &
{\tabfont Involuntary errors due to inattention or distraction} &
{\tabfont Improve workspace, process design, and accountability}\\
\rowcolor{paleblue}
{\tabfont \textbf{2a) Mistakes --- rule-based}} &
{\tabfont Voluntary errors arising from flawed rules} &
{\tabfont Enhance rules and controls}\\
\rowcolor{paleblue}
{\tabfont \textbf{2b) Mistakes --- knowledge-based}} &
{\tabfont Errors in new situations due to wrong actions} &
{\tabfont Better documentation, training, and escalation processes}\\
{\tabfont \textbf{3) Violations}} &
{\tabfont Deliberate deviations from rules} &
{\tabfont Strengthen supervisory controls and detection methods}\\
\end{tabularx}
\end{minipage}
\orfigcap{Slips are involuntary, mistakes are voluntary, violations are deliberate.}

\orsub{Process improvement techniques}

\noindent\begin{minipage}{\textwidth}
\renewcommand{\arraystretch}{1.35}
\begin{tabularx}{\textwidth}{@{}L{0.22\textwidth}YL{0.26\textwidth}@{}}
\rowcolor{navy} \thd{Method} & \thd{What it does} & \thd{Best for}\\
{\tabfont \textbf{1) Lean Six Sigma (DMAIC)}} &
{\tabfont Structured approach for complex processes: eliminate waste (Lean) and
minimise variation (Six Sigma). DMAIC cycle: Define, Measure, Analyse, Improve,
Control.} &
{\tabfont Manufacturing, complex product development}\\
\rowcolor{paleblue}
{\tabfont \textbf{2) Quality Improvement (PDSA)}} &
{\tabfont Flexible cycle for continuous improvement, through ongoing cycles.
PDSA cycle: Plan, Do, Study, Act.} &
{\tabfont Various industries, adaptable to different situations}\\
\end{tabularx}
\end{minipage}
\orfigcap{DMAIC for structure, PDSA for flexibility.}

% ---------------------------------------------------------------------
\lohead{44}{e}{Explain how operational risk can arise with new products, new
business initiatives, or mergers and acquisitions, and describe ways to mitigate
these risks.}

\orsub{Assessing new products and initiatives}

\begin{orkeybox}
{\sffamily\bfseries\fontsize{8.6}{10}\selectfont\color{navy}NPAP and NIRAP}\par
\vspace{2pt}
A New Initiative Risk Assessment Process (NIRAP) business case outlines the
purpose, alternatives, benefits, costs, and risks, along with mitigation
strategies, for new products or initiatives. This sits alongside the New Product
Approval Process (NPAP).
\end{orkeybox}

New products and initiatives can change existing risks or introduce new ones, so
ORM teams help identify, assess, and mitigate these risks.

\begin{lettered}
\item \textbf{Initial stage:} identify and assess risks, create mitigation plans.
\item \textbf{Project life:} regular reporting, meetings with risk and project
teams.
\item \textbf{Project closure:} debriefs, evaluations, risk analysis, and lessons
learned.
\end{lettered}

\orsub{Mergers and acquisitions (M\&A)}

\begin{lettered}
\item Mergers and acquisitions entail inheriting credit, market, and operational
risks from the acquired firm, posing challenges for risk assessment post-merger.
\item Integrating the acquired firm can introduce operational risks such as errors
in system integration and communication breakdowns, necessitating the role of
operational risk management in risk identification and mitigation.
\end{lettered}

% ---------------------------------------------------------------------
\lohead{44}{f}{Identify and describe approaches firms should use to mitigate the
impact of operational risk events.}

Financial firms know that even strong controls can't prevent all losses. Here is
how they can reduce the impact of operational risk events.

\orsub{1) Contingency planning}

\begin{ordefbox}{Business Continuity Management (BCM)}
A firm's action plan for crises. It identifies critical business areas, develops a
strategy, and implements a plan.
\end{ordefbox}

\begin{lettered}
\item Involves senior management, a designated owner, and risk identification
(technological, reputational, environmental).
\item Develops strategies and plans, and implements them.
\item Additional mitigation options: external insurance and crisis management
specialists.
\end{lettered}

\orsub{2) Event and crisis management}

\orsubb{Three keys to a strong Business Continuity Plan (BCP)}
\begin{lettered}
\item \textbf{Speed:} react quickly to crises.
\item \textbf{Competent teams:} assign tasks to specialists.
\item \textbf{Transparency:} firms should promptly and transparently disclose
operational loss events to staff, external stakeholders, and the public.
\end{lettered}

\orsubb{Team structure for crisis management}
\begin{lettered}
\item \textbf{Technical team:} handles IT-related disruptions; includes IT and
security specialists.
\item \textbf{Communications team:} manages media and stakeholder communication,
reporting to senior leadership.
\end{lettered}

\orsubb{Phases of crisis management}

\begin{center}
\begin{tikzpicture}[font=\fontsize{7.8}{9.6}\selectfont,
  st/.style={text width=3.2cm,align=center,rounded corners=3pt,inner sep=5pt,
             text=white,font=\sffamily\bfseries\fontsize{8.2}{10}\selectfont},
  sb/.style={text width=3.2cm,align=center,rounded corners=3pt,inner sep=6pt,
             fill=paleblue,draw=palenavy,line width=0.6pt,text=navy}]
\foreach \i/\x/\c/\t/\d in {%
  1/-6.0/navy/{a) Crisis}/{Initial recognition of the incident},
  2/-2.0/teal/{b) Emergency response}/{Specialists act to mitigate the crisis},
  3/2.0/forest/{c) Recovery}/{Operations begin to resume, guided by RPO and RTO},
  4/6.0/olive/{d) Restoration}/{Full restoration of functions, return to normal operations}}
{
  \node[st,fill=\c,anchor=north] (t\i) at (\x,0) {\t};
  \node[sb,anchor=north] (b\i) at (t\i.south) {\d};
}
\foreach \a/\b in {1/2,2/3,3/4}
  \draw[-{Stealth[length=4pt]},navy!35,line width=1.2pt]
    ($(t\a.east)+(0.04,0)$) -- ($(t\b.west)+(-0.04,0)$);
\end{tikzpicture}
\end{center}
\orfigcap{The four phases of crisis management. RPO is the Recovery Point
Objective; RTO is the Recovery Time Objective.}

% ---------------------------------------------------------------------
\lohead{44}{g}{Describe methods for the transfer of operational risks and the
management of reputational risk, and assess their effectiveness in different
situations.}

\orsub{Risk transfer}

\noindent\begin{minipage}{\textwidth}
\renewcommand{\arraystretch}{1.32}
\begin{tabularx}{\textwidth}{@{}L{0.18\textwidth}YY@{}}
\rowcolor{navy} \thd{Method} & \thd{Advantage} & \thd{Disadvantages}\\
{\tabfont \textbf{1) External insurance}} &
{\tabfont Reduced loss volatility, via premiums paid for financial
compensation.} &
{\tabfont
a)~May not fully cover losses.\newline
b)~Delays in compensation can cause cash flow issues.\newline
c)~Insurance may not be cost-effective for all risks.\newline
d)~Not all risks are fully transferable.}\\
\rowcolor{paleblue}
{\tabfont \textbf{2) Outsourcing}} &
{\tabfont Access to specialised skills and potentially lower costs.} &
{\tabfont
a)~New risks arise from relying on a third party (e.g., service delays, security
breaches).\newline
b)~Outsourcing may be risk sharing, not full transfer.\newline
c)~Risks transferred depend on the task and the third party's control
effectiveness.}\\
\end{tabularx}
\end{minipage}
\orfigcap{Insurance and outsourcing as transfer mechanisms.}

\orsub{Managing reputational risk}

\orsubb{1) Preventive strategies}
\begin{lettered}
\item Build customer confidence and ensure transparency.
\item Integrate reputation management into employee roles and rewards.
\item Conduct detailed due diligence and analysis before forming new business
relationships.
\end{lettered}

\orsubb{2) Crisis communication}
\begin{lettered}
\item Apply the three Rs: \textbf{Regret}, \textbf{Reason}, and \textbf{Remedy}.
\item Use targeted communication based on stakeholder analysis.
\end{lettered}

\orsubb{3) Long-term relationship and transparency}
\begin{lettered}
\item Maintain strong relationships with all stakeholders, including regulators.
\item Use reputational capital built over time to enhance resilience and
potentially improve reputation post-crisis.
\end{lettered}
